% Part: lambda-calculus
% Chapter: introduction
% Section: free-variables

\documentclass[../../../include/open-logic-section]{subfiles}

\begin{document}

\olfileid[hi]{lam}{syn}{fv}
\olsection{मुक्त चर}

लैम्ब्डा कलन फलनों के विषय में है, और लैम्ब्डा अमूर्तन से ही फलन उत्पन्न
होते हैं। सहज रूप से, $\lambd[x][M]$ वह फलन है जिसके मान~$M$ द्वारा तब दिए
जाते हैं जब फलन का तर्क~$x$ को नियत किया जाए। किंतु~$M$ में~$x$ का प्रत्येक
आविर्भाव प्रासंगिक नहीं है: यदि $M$ में कोई दूसरा अमूर्त
$\lambd[x][N]$ हो, तो~$N$ में $x$ के आविर्भाव $\lambd[x][N]$ के लिए
प्रासंगिक हैं, $\lambd[x][M]$ के लिए नहीं। अतः $\lambd[x][M]$ के भीतर का
लैम्ब्डा अमूर्त $\lambd[x]$, ~$M$ में $x$ के उन आविर्भावों को \emph{बाँधता}
है जो पहले से किसी अन्य लैम्ब्डा अमूर्त द्वारा आबद्ध नहीं हैं---अर्थात्
~$M$ में~$x$ के \emph{मुक्त} आविर्भावों को।

\begin{defn}[व्याप्ति]
यदि $\lambd[x][M]$ किसी पद~$N$ के भीतर आता है, तो उस आविर्भाव के भीतर का
~$M$, उस~$\lambd[x]$ की \emph{व्याप्ति} है।
\end{defn}


\begin{defn}[मुक्त और आबद्ध आविर्भाव]
  किसी पद~$M$ में चर $x$ का कोई आविर्भाव \emph{मुक्त} है यदि वह किसी
  $\lambd[x]$ की व्याप्ति में नहीं है; अन्यथा वह \emph{आबद्ध} है।
  $\lambd[x][M]$ में चर~$x$ का कोई आविर्भाव आरंभिक $\lambd[x]$ द्वारा तभी
  और केवल तभी आबद्ध है जब~$M$ में $x$ का वह आविर्भाव मुक्त हो।
\end{defn}

\begin{ex}
  $\lambd[x][x y]$ में $x$ और $y$ दोनों $\lambd[x]$ की व्याप्ति में हैं,
  इसलिए $x$, ~$\lambd[x]$ द्वारा आबद्ध है। चूँकि $y$ किसी $\lambd[y]$ की
  व्याप्ति में नहीं है, इसलिए वह मुक्त है। $\lambd[x][x x]$ में $x$ के
  दोनों आविर्भाव~$\lambd[x]$ द्वारा आबद्ध हैं, क्योंकि $xx$ में दोनों मुक्त
  हैं। $((\lambd[x][xx])x)$ में~$x$ का अंतिम आविर्भाव मुक्त है, क्योंकि वह
  किसी $\lambd[x]$ की व्याप्ति में नहीं है।
  $\lambd[x][(\lambd[x][x])x]$ में पहले $\lambd[x]$ की व्याप्ति
  $(\lambd[x][x])x$ है और दूसरे $\lambd[x]$ की व्याप्ति~$x$ का अंतिम से
  पहले वाला आविर्भाव है। $(\lambd[x][x])x$ में $x$ का अंतिम आविर्भाव मुक्त
  और अंतिम से पहले वाला आबद्ध है। अतः $\lambd[x][(\lambd[x][x])x]$ में
  $x$ का अंतिम से पहले वाला आविर्भाव दूसरे $\lambd[x]$ द्वारा और अंतिम
  आविर्भाव पहले $\lambd[x]$ द्वारा आबद्ध है।
\end{ex}

किसी पद $P$ में चरों के सभी आविर्भाव जाँचकर हम मुक्त चरों का समुच्चय पा
सकते हैं। इस समुच्चय को $\FV{P}$ से दर्शाते हैं; उसकी स्वाभाविक परिभाषा
इस प्रकार है:

\begin{defn}[किसी पद के मुक्त चर] \ollabel{def:fv}
  किसी पद के \emph{मुक्त चरों} का समुच्चय आगमनतः इस प्रकार परिभाषित है:
  \begin{enumerate}
    \item $\FV{x} = \{x\}$ \ollabel{def:fv1} 
    \item $\FV{\lambd[x][N]} = \FV{N} \setminus \{x\}$    \ollabel{def:fv2}
    \item $\FV{PQ} = \FV{P} \cup \FV{Q}$ \ollabel{def:fv3}
  \end{enumerate}
\end{defn}

\begin{prob}
  \begin{enumerate}
  \item इस पद में $\lambd[g]$ और दोनों $\lambd[x]$ की व्याप्तियाँ पहचानें:
    $\lambd[g][(\lambd[x][g (x x)]) \lambd[x][g (x x)]]$।
  \item $\lambd[g][(\lambd[x][g (x x)]) \lambd[x][g (x x)]]$ में क्या
    चरों के सभी आविर्भाव आबद्ध हैं? वे क्रमशः किन अमूर्तनों द्वारा आबद्ध हैं?
    \item $\FV{\lambd[x][(\lambd[y][(\lambd[z][x y]) z]) y]}$ दें।
  \end{enumerate}
\end{prob}

\begin{explain}
कोई मुक्त चर बाहरी संसार (\emph{परिवेश}) के संदर्भ जैसा है, और मुक्त चर
वाले पद को आंशिक रूप से विनिर्दिष्ट पद माना जा सकता है, क्योंकि उसका
व्यवहार इस पर निर्भर है कि हम परिवेश कैसे बनाते हैं। उदाहरणार्थ, पद
$\lambd[x][f x]$ में चर~$f$ मुक्त है; यह पद तर्क~$x$ स्वीकार करके उस तर्क
पर $f$ का मान लौटाता है। इस पद का मान उस परिवेश पर निर्भर है जिसमें वह है,
विशेषतः उस परिवेश में $f$ के मान पर।

यदि हम इस पद पर अमूर्तन लागू करें, तो $\lambd[f][\lambd[x][f x]]$ मिलता
है। अब यह पद परिवेश-चर $f$ पर निर्भर नहीं है, क्योंकि अब यह ऐसे फलन को
दर्शाता है जो दो तर्क स्वीकार करता और पहले को दूसरे पर लागू करने का परिणाम
लौटाता है। परिवेश में $f$ को बदलने से इस पद के व्यवहार पर कोई प्रभाव नहीं
होगा, क्योंकि पद केवल तर्क के रूप में दी गई वस्तु का उपयोग करेगा, परिवेश
में $f$ के मान का नहीं।
\end{explain}

\begin{defn}[संवृत पद, संयोजक]
  जिस पद में कोई मुक्त चर न हो उसे \emph{संवृत पद}, अथवा
  \emph{संयोजक}, कहते हैं।
\end{defn}

\begin{lem}\ollabel{lem:fv}
  \begin{enumerate}
    \item \ollabel{lem:fv-abs} यदि $y \neq x$, तो $y \in
      \FV{\lambd[x][N]}$ तभी और केवल तभी जब $y \in \FV{N}$।
    \item \ollabel{lem:fv-app} $y \in \FV{PQ}$ तभी और केवल तभी जब $y \in
      \FV{P}$ अथवा $y \in \FV{Q}$।
    \end{enumerate}
\end{lem}

\begin{proof}
  अभ्यास।
\end{proof}

\begin{prob}
  \olref[lam][syn][fv]{lem:fv} सिद्ध करें।
\end{prob}

\end{document}
